% GENERATED FILE. Edit canonical lessons, then run npm run generate:books.
% canonical-source-hash: fnv1a64:91f317bb93bcb71e
% canonical-lessons: MR-C380-anivarya, MR-C380-aicchik, MR-C380-apariharya, MR-C380-ladne, MR-C380-tartud

\chapter{Must or May}
\label{ch:mr-must-or-may}
\hlchaptermodality{380}

\begin{chapteropening}
\textbf{By the end of this chapter:} \emph{I can say what is compulsory, optional or unavoidable.}

Complete the last lesson of chapter 380: I can say what is compulsory, optional or unavoidable.
\end{chapteropening}

\section[\texorpdfstring{\mr{अनिवार्य} (anivārya)}{अनिवार्य (anivārya)}]{\mr{अनिवार्य} (anivārya) — compulsory}

\label{lesson:MR-C380-anivarya}



\begin{quote}
Before the new one: say the Marathi for to relax, then the Marathi for to be manageable.
\end{quote}

\subsection*{You'll want to know: \mr{अनिवार्य}}
\textbf{\mr{अनिवार्य}} — \emph{anivārya} — \textquotedblleft{}compulsory\textquotedblright{}.

\begin{quote}
\mr{शाळेत} \mr{मराठी} \mr{हा} \mr{विषय} \mr{अनिवार्य} \mr{आहे}.
\end{quote}

\begin{quote}
\emph{śāḷet marāṭhī hā viṣay anivārya āhe.}
\end{quote}

\begin{quote}
Marathi is a compulsory subject at school.
\end{quote}

Invariable; in speech \mr{सक्तीचं}.

\begin{cousinweb}
\textbf{\mr{अनिवार्य}} (\emph{anivārya}) is a tatsama, taken from Sanskrit \emph{anivārya}, \textquotedblleft{}not to be kept off, unavoidable\textquotedblright{} (Monier-Williams).
\end{cousinweb}

\begin{grammarlens}[title={What you've built}]
One more word for saying what happens if.
\end{grammarlens}

\subsection*{Your turn}
\begin{itemize}
\raggedright
  \item \emph{Say it:} \emph{anivārya}
  \item \emph{Say it:} \emph{anivārya}, once more
  \item \emph{Say it:} \emph{śāḷet marāṭhī hā viṣay anivārya āhe.} — Marathi is a compulsory subject at school.
  \item \emph{Recall:} say the Marathi for to relax, then the Marathi for to be manageable, then say \emph{anivārya} again
\end{itemize}

\begin{tcolorbox}[breakable,colback=teal!4,colframe=teal!35!black,title={Before you move on}]
What does \mr{अनिवार्य} mean? (\textquotedblleft{}Compulsory\textquotedblright{}.) Say it once more.
\end{tcolorbox}
\section[\texorpdfstring{\mr{ऐच्छिक} (aicchik)}{ऐच्छिक (aicchik)}]{\mr{ऐच्छिक} (aicchik) — optional}

\label{lesson:MR-C380-aicchik}



\begin{quote}
Before the new one: say the Marathi for to be manageable, then the Marathi for compulsory.
\end{quote}

\subsection*{You'll want to know: \mr{ऐच्छिक}}
\textbf{\mr{ऐच्छिक}} — \emph{aicchik} — \textquotedblleft{}optional\textquotedblright{}.

\begin{quote}
\mr{हा} \mr{विषय} \mr{ऐच्छिक} \mr{आहे}, \mr{तो} \mr{घेतला} \mr{नाही} \mr{तरी} \mr{चालेल}.
\end{quote}

\begin{quote}
\emph{hā viṣay aicchik āhe, to ghetlā nāhī tarī cālel.}
\end{quote}

\begin{quote}
This subject is optional; it is fine if you don't take it.
\end{quote}

Built on \mr{इच्छा}, wish: by choice.

\begin{cousinweb}
\textbf{\mr{ऐच्छिक}} (\emph{aicchik}) is a tatsama, taken from Sanskrit \emph{aicchika}, \textquotedblleft{}optional, at will\textquotedblright{} (Monier-Williams). It is built on \emph{icchā}, \textquotedblleft{}wish\textquotedblright{}.
\end{cousinweb}

\begin{grammarlens}[title={What you've built}]
One more word for saying what happens if.
\end{grammarlens}

\subsection*{Your turn}
\begin{itemize}
\raggedright
  \item \emph{Say it:} \emph{aicchik}
  \item \emph{Say it:} \emph{aicchik}, once more
  \item \emph{Say it:} \emph{hā viṣay aicchik āhe, to ghetlā nāhī tarī cālel.} — This subject is optional; it is fine if you don't take it.
  \item \emph{Recall:} say the Marathi for to be manageable, then the Marathi for compulsory, then say \emph{aicchik} again
\end{itemize}

\begin{tcolorbox}[breakable,colback=teal!4,colframe=teal!35!black,title={Before you move on}]
What does \mr{ऐच्छिक} mean? (\textquotedblleft{}Optional\textquotedblright{}.) Say it once more.
\end{tcolorbox}
\section[\texorpdfstring{\mr{अपरिहार्य} (aparihārya)}{अपरिहार्य (aparihārya)}]{\mr{अपरिहार्य} (aparihārya) — unavoidable}

\label{lesson:MR-C380-apariharya}



\begin{quote}
Before the new one: say the Marathi for compulsory, then the Marathi for optional.
\end{quote}

\subsection*{You'll want to know: \mr{अपरिहार्य}}
\textbf{\mr{अपरिहार्य}} — \emph{aparihārya} — \textquotedblleft{}unavoidable\textquotedblright{}.

\begin{quote}
\mr{पावसामुळे} \mr{थोडा} \mr{उशीर} \mr{अपरिहार्य} \mr{होता}.
\end{quote}

\begin{quote}
\emph{pāvsāmuḷe thoḍā uśīr aparihārya hotā.}
\end{quote}

\begin{quote}
Because of the rain, a little delay was unavoidable.
\end{quote}

Invariable; \mr{अ} + \mr{परिहार्य}, not avoidable.

\begin{cousinweb}
\textbf{\mr{अपरिहार्य}} (\emph{aparihārya}) is a tatsama, taken from Sanskrit \emph{aparihārya}, \textquotedblleft{}not to be avoided\textquotedblright{} (Monier-Williams).
\end{cousinweb}

\begin{grammarlens}[title={What you've built}]
One more word for saying what happens if.
\end{grammarlens}

\subsection*{Your turn}
\begin{itemize}
\raggedright
  \item \emph{Say it:} \emph{aparihārya}
  \item \emph{Say it:} \emph{aparihārya}, once more
  \item \emph{Say it:} \emph{pāvsāmuḷe thoḍā uśīr aparihārya hotā.} — Because of the rain, a little delay was unavoidable.
  \item \emph{Recall:} say the Marathi for compulsory, then the Marathi for optional, then say \emph{aparihārya} again
\end{itemize}

\begin{tcolorbox}[breakable,colback=teal!4,colframe=teal!35!black,title={Before you move on}]
What does \mr{अपरिहार्य} mean? (\textquotedblleft{}Unavoidable\textquotedblright{}.) Say it once more.
\end{tcolorbox}
\section[\texorpdfstring{\mr{लादणे} (lādṇe)}{लादणे (lādṇe)}]{\mr{लादणे} (lādṇe) — to impose}

\label{lesson:MR-C380-ladne}



\begin{quote}
Before the new one: say the Marathi for optional, then the Marathi for unavoidable.
\end{quote}

\subsection*{You'll want to know: \mr{लादणे}}
\textbf{\mr{लादणे}} — \emph{lādṇe} — \textquotedblleft{}to impose\textquotedblright{}.

\begin{quote}
\mr{आईवडिलांनी} \mr{मुलांवर} \mr{आपली} \mr{मतं} \mr{लादू} \mr{नयेत}.
\end{quote}

\begin{quote}
\emph{āīvḍilānnī mulāṁvar āplī mataṁ lādū nayet.}
\end{quote}

\begin{quote}
Parents should not impose their views on their children.
\end{quote}

Past, with -\mr{ने} (he …): \mr{त्याने} \mr{लादलं}, \emph{tyāne lādlaṁ}.

The one burdened takes -\mr{वर}; also to load a cart.

\begin{cousinweb}
\textbf{\mr{लादणे}} (\emph{lādṇe}) is a tadbhava: CDIAL lists Marathi \emph{lādṇẽ}, \textquotedblleft{}to heap up, load\textquotedblright{}, under Sanskrit \emph{lardayati}, \textquotedblleft{}loads\textquotedblright{} (CDIAL 10966).
\end{cousinweb}

\begin{grammarlens}[title={What you've built}]
One more word for saying what happens if.
\end{grammarlens}

\subsection*{Your turn}
\begin{itemize}
\raggedright
  \item \emph{Say it:} \emph{lādṇe}
  \item \emph{Say it:} \emph{lādṇe}, once more
  \item \emph{Say it:} \emph{āīvḍilānnī mulāṁvar āplī mataṁ lādū nayet.} — Parents should not impose their views on their children.
  \item \emph{Recall:} say the Marathi for optional, then the Marathi for unavoidable, then say \emph{lādṇe} again
\end{itemize}

\begin{tcolorbox}[breakable,colback=teal!4,colframe=teal!35!black,title={Before you move on}]
What does \mr{लादणे} mean? (\textquotedblleft{}To impose\textquotedblright{}.) Say it once more.
\end{tcolorbox}
\section[\texorpdfstring{\mr{तरतूद} (tartūd)}{तरतूद (tartūd)}]{\mr{तरतूद} (tartūd) — a provision (in a law or budget)}

\label{lesson:MR-C380-tartud}



\begin{quote}
Before the new one: say the Marathi for unavoidable, then the Marathi for to impose.
\end{quote}

\subsection*{You'll want to know: \mr{तरतूद}}
\textbf{\mr{तरतूद}} — \emph{tartūd} — \textquotedblleft{}a provision (in a law or budget)\textquotedblright{}.

\begin{quote}
\mr{अर्थसंकल्पात} \mr{शिक्षणासाठी} \mr{मोठी} \mr{तरतूद} \mr{आहे}.
\end{quote}

\begin{quote}
\emph{arthasaṅkalpāt śikṣaṇāsāṭhī moṭhī tartūd āhe.}
\end{quote}

\begin{quote}
The budget has a large provision for education.
\end{quote}

Feminine.

X \mr{ची} \mr{तरतूद} \mr{करणे}, to make provision for X.

\begin{cousinweb}
\textbf{\mr{तरतूद}} (\emph{tartūd}) is a Marathi noun whose origin is not settled.
\end{cousinweb}

\begin{grammarlens}[title={What you've built}]
One more word for saying what happens if.
\end{grammarlens}

\subsection*{Your turn}
\begin{itemize}
\raggedright
  \item \emph{Say it:} \emph{tartūd}
  \item \emph{Say it:} \emph{tartūd}, once more
  \item \emph{Say it:} \emph{arthasaṅkalpāt śikṣaṇāsāṭhī moṭhī tartūd āhe.} — The budget has a large provision for education.
  \item \emph{Recall:} say the Marathi for unavoidable, then the Marathi for to impose, then say \emph{tartūd} again
\end{itemize}

\begin{tcolorbox}[breakable,colback=teal!4,colframe=teal!35!black,title={Before you move on}]
What does \mr{तरतूद} mean? (\textquotedblleft{}A provision (in a law or budget)\textquotedblright{}.) Say it once more.
\end{tcolorbox}
